\documentclass{article}
\usepackage{amsmath,amssymb,amsthm}
\usepackage{algorithm,algorithmic}
\usepackage{enumitem}
\newtheorem{theorem}{Theorem}
\newtheorem{lemma}{Lemma}
\newtheorem{assumption}{Assumption}
\newtheorem{remark}{Remark}
\begin{document}

\section*{Proximal Gradient Method}

The method solves the composite convex optimization problem  

\[
\min_{x\in\mathbb{R}^d}\; \Phi(x)\;:=\;f(x)+g(x),
\tag{P}
\]

where  

\begin{itemize}
  \item $f:\mathbb{R}^d\to\mathbb{R}$ is convex and differentiable with $L$‑Lipschitz continuous gradient,
  \item $g:\mathbb{R}^d\to\mathbb{R}\cup\{+\infty\}$ is proper, closed and convex (possibly nonsmooth).
\end{itemize}

The algorithm alternates a gradient step on $f$ with a proximal step on $g$.

\section*{Mathematical Formulation}

Given a stepsize $\eta>0$, the iteration is  

\[
x_{k+1}\;=\;\operatorname{prox}_{\eta g}\!\bigl(x_k-\eta\nabla f(x_k)\bigr),\qquad k=0,1,2,\dots
\tag{1}
\]

The proximal operator is defined for any $z\in\mathbb{R}^d$ by  

\[
\operatorname{prox}_{\eta g}(z)\;:=\;\arg\min_{u\in\mathbb{R}^d}\Bigl\{ g(u)+\frac{1}{2\eta}\|u-z\|^2\Bigr\}.
\tag{2}
\]

\section*{Pseudocode}

\begin{algorithm}[H]
\caption{Proximal Gradient Method}
\begin{algorithmic}[1]
\STATE \textbf{Input:} stepsize $\eta>0$, initial point $x_0\in\mathbb{R}^d$, maximum iterations $K$.
\FOR{$k=0$ \textbf{to} $K-1$}
\STATE Compute gradient $g_k\gets\nabla f(x_k)$.
\STATE Gradient step $y_k\gets x_k-\eta g_k$.
\STATE Proximal step $x_{k+1}\gets\operatorname{prox}_{\eta g}(y_k)$.
\ENDFOR
\STATE \textbf{Output:} $x_K$ (or the best iterate).
\end{algorithmic}
\end{algorithm}

\section*{Dry Run Example}

Consider the scalar problem  

\[
f(w)=(w-3)^2,\qquad g\equiv 0,\qquad \eta=0.1,\qquad w_0=0.
\]

Since $g\equiv0$, $\operatorname{prox}_{\eta g}$ is the identity.  
The gradient $\nabla f(w)=2(w-3)$.

\begin{center}
\begin{tabular}{c|c|c|c}
Iteration $k$ & $w_k$ & $\nabla f(w_k)$ & $w_{k+1}=w_k-\eta\nabla f(w_k)$\\ \hline
0 & $0$ & $-6$ & $0-0.1(-6)=0.6$ \\
1 & $0.6$ & $-4.8$ & $0.6-0.1(-4.8)=1.08$ \\
2 & $1.08$ & $-3.84$ & $1.08-0.1(-3.84)=1.464$ \\
3 & $1.464$ & $-3.072$ & $1.464-0.1(-3.072)=1.7712$
\end{tabular}
\end{center}

Objective values $\Phi(w_k)= (w_k-3)^2$ decrease from $9$ to $2.944$ after three steps, illustrating convergence toward the minimiser $w^\star=3$.

\section*{Assumptions}

\begin{assumption}[Smoothness of $f$]\label{ass:smooth}
$f$ is convex and differentiable with $L$‑Lipschitz gradient, i.e.  

\[
\|\nabla f(x)-\nabla f(y)\|\le L\|x-y\|,\qquad\forall x,y\in\mathbb{R}^d.
\tag{A1}
\]

\end{assumption}

\begin{assumption}[Convexity of $g$]\label{ass:convex}
$g$ is proper, closed and convex. No smoothness is required.
\tag{A2}
\end{assumption}

\begin{assumption}[Stepsize]\label{ass:stepsize}
The constant stepsize satisfies  

\[
0<\eta\le \frac{1}{L}.
\tag{A3}
\]

\end{assumption}

\section*{Main Convergence Theorem}

\begin{theorem}\label{thm:conv}
Let $\{x_k\}$ be generated by \eqref{1} under Assumptions \ref{ass:smooth}–\ref{ass:stepsize}. Denote $x^\star\in\arg\min\Phi$ and $\Phi^\star:=\Phi(x^\star)$. Then for all $K\ge1$

\[
\Phi(x_K)-\Phi^\star\;\le\;\frac{\|x_0-x^\star\|^2}{2\eta K}.
\tag{3}
\]

Consequently $\Phi(x_K)\to\Phi^\star$ at the rate $O(1/K)$.
\end{theorem}

\section*{Theoretical Properties}

\begin{itemize}
  \item \textbf{Stability.} The iteration is non‑expansive in the sense that  

  \[
  \|x_{k+1}-x^\star\|\le\|x_k-x^\star\|,
  \]
  which follows directly from the non‑expansiveness of the proximal mapping (Lemma~\ref{lem:nonexpansive}) and the $L$‑smoothness of $f$.
  \item \textbf{Hyperparameter Sensitivity.} The bound \eqref{3} requires $\eta\le1/L$. Choosing $\eta$ close to $1/L$ yields the smallest constant $\frac{1}{2\eta}$, i.e. the fastest theoretical decay.
  \item \textbf{Comparison to Gradient Descent.} When $g\equiv0$, \eqref{1} reduces to standard gradient descent and the same $O(1/K)$ rate is recovered. For nonsmooth $g$, the proximal step allows us to retain this rate without requiring subgradient steps.
\end{itemize}

\section*{Discussion}

The theorem shows that the proximal gradient method enjoys the same sublinear convergence as vanilla gradient descent for convex smooth problems, despite handling a nonsmooth component $g$. The proof hinges on two elementary tools: the descent lemma for $L$‑smooth functions and the non‑expansiveness of the proximal operator. Extensions to accelerated rates ($O(1/K^2)$) or stochastic settings follow similar patterns but require additional momentum or variance‑reduction arguments.

\section*{Preliminaries (Definitions and Lemmas)}

\begin{definition}[Proximal Operator]\label{def:prox}
For $\eta>0$ and convex $g$,  

\[
\operatorname{prox}_{\eta g}(z)=\arg\min_{u}\Bigl\{g(u)+\frac{1}{2\eta}\|u-z\|^2\Bigr\}.
\]
\end{definition}

\begin{lemma}[Non‑expansiveness of $\operatorname{prox}_{\eta g}$]\label{lem:nonexpansive}
For any $z_1,z_2\in\mathbb{R}^d$,
\[
\bigl\|\operatorname{prox}_{\eta g}(z_1)-\operatorname{prox}_{\eta g}(z_2)\bigr\|
\le\|z_1-z_2\|.
\tag{4}
\]
\end{lemma}
\begin{proof}
Let $u_i=\operatorname{prox}_{\eta g}(z_i)$, $i=1,2$. By the optimality condition,
\[
0\in\partial g(u_i)+\frac{1}{\eta}(u_i-z_i).
\]
Thus $\frac{1}{\eta}(z_i-u_i)\in\partial g(u_i)$. Using monotonicity of $\partial g$,
\[
\big\langle \frac{1}{\eta}(z_1-u_1)-\frac{1}{\eta}(z_2-u_2),\,u_1-u_2\big\rangle\ge0.
\]
Rearranging yields \eqref{4}. ∎
\end{proof}

\begin{lemma}[Descent Lemma for $L$‑smooth $f$]\label{lem:descent}
For any $x,y\in\mathbb{R}^d$,
\[
f(y)\le f(x)+\langle\nabla f(x),y-x\rangle+\frac{L}{2}\|y-x\|^2.
\tag{5}
\]
\end{lemma}
\begin{proof}
Standard consequence of $L$‑Lipschitz gradient; see e.g. Nesterov (2004). ∎
\end{proof}

\section*{Main Proof (Step‑by‑Step Derivation)}

Let $x^\star\in\arg\min\Phi$ and denote $y_k:=x_k-\eta\nabla f(x_k)$.  
From the definition \eqref{1}, $x_{k+1}= \operatorname{prox}_{\eta g}(y_k)$.  

\begin{enumerate}[label={[\textbf{Step \arabic*}]}]
\item \textbf{Apply non‑expansiveness.} Using Lemma~\ref{lem:nonexpansive} with $z_1=y_k$, $z_2=x^\star-\eta\nabla f(x^\star)$ (note that $x^\star$ satisfies the optimality condition $0\in\nabla f(x^\star)+\partial g(x^\star)$, hence $x^\star=\operatorname{prox}_{\eta g}(x^\star-\eta\nabla f(x^\star))$), we obtain  

\[
\|x_{k+1}-x^\star\|^2
\le\|y_k-(x^\star-\eta\nabla f(x^\star))\|^2.
\tag{6}
\]

\item \textbf{Expand the right‑hand side.}  
\[
\begin{aligned}
\|y_k-(x^\star-\eta\nabla f(x^\star))\|^2
&=\|x_k-\eta\nabla f(x_k)-x^\star+\eta\nabla f(x^\star)\|^2\\
&=\|x_k-x^\star\|^2-2\eta\langle\nabla f(x_k)-\nabla f(x^\star),x_k-x^\star\rangle\\
&\qquad+\eta^2\|\nabla f(x_k)-\nabla f(x^\star)\|^2.
\end{aligned}
\tag{7}
\]

\item \textbf{Use $L$‑smoothness.} By the $L$‑Lipschitz property,  

\[
\|\nabla f(x_k)-\nabla f(x^\star)\|^2\le L^2\|x_k-x^\star\|^2.
\tag{8}
\]

\item \textbf{Combine (6)–(8).} Substituting (7) and (8) into (6) yields  

\[
\|x_{k+1}-x^\star\|^2
\le\|x_k-x^\star\|^2-2\eta\langle\nabla f(x_k)-\nabla f(x^\star),x_k-x^\star\rangle
+\eta^2L^2\|x_k-x^\star\|^2.
\tag{9}
\]

\item \textbf{Rearrange terms.} Because $\eta\le 1/L$ (Assumption \ref{ass:stepsize}), we have $\eta^2L^2\le\eta L$. Hence  

\[
\|x_{k+1}-x^\star\|^2
\le\|x_k-x^\star\|^2-2\eta\langle\nabla f(x_k)-\nabla f(x^\star),x_k-x^\star\rangle
+\eta L\|x_k-x^\star\|^2.
\tag{10}
\]

\item \textbf{Insert convexity of $f$.} Convexity gives  

\[
\langle\nabla f(x_k)-\nabla f(x^\star),x_k-x^\star\rangle\ge f(x_k)-f(x^\star).
\tag{11}
\]

\item \textbf{Insert convexity of $g$.} The optimality condition for $x^\star$ implies  

\[
g(x_{k+1})-g(x^\star)\le\langle s^\star,x_{k+1}-x^\star\rangle,
\qquad s^\star\in\partial g(x^\star)=-\nabla f(x^\star).
\tag{12}
\]

Adding (11) and (12) yields  

\[
\Phi(x_k)-\Phi^\star\le\langle\nabla f(x_k)+s^\star,\,x_k-x^\star\rangle.
\tag{13}
\]

\item \textbf{Bound the decrease of the distance.} Using (10) together with (13),

\[
\begin{aligned}
\|x_{k+1}-x^\star\|^2
&\le\|x_k-x^\star\|^2-2\eta\bigl(\Phi(x_k)-\Phi^\star\bigr)+\eta L\|x_k-x^\star\|^2\\
&=\bigl(1+\eta L\bigr)\|x_k-x^\star\|^2-2\eta\bigl(\Phi(x_k)-\Phi^\star\bigr).
\end{aligned}
\tag{14}
\]

\item \textbf{Use $\eta\le 1/L$ to simplify.} Since $1+\eta L\le 2$, we obtain  

\[
\|x_{k+1}-x^\star\|^2
\le\|x_k-x^\star\|^2-2\eta\bigl(\Phi(x_k)-\Phi^\star\bigr).
\tag{15}
\]

\item \textbf{Telescoping sum.} Summing \eqref{15} for $k=0,\dots,K-1$ gives  

\[
\|x_{K}-x^\star\|^2
\le\|x_{0}-x^\star\|^2-2\eta\sum_{k=0}^{K-1}\bigl(\Phi(x_k)-\Phi^\star\bigr).
\tag{16}
\]

Discard the non‑negative term $\|x_{K}-x^\star\|^2$ and divide by $2\eta K$:

\[
\frac{1}{K}\sum_{k=0}^{K-1}\bigl(\Phi(x_k)-\Phi^\star\bigr)
\le\frac{\|x_{0}-x^\star\|^2}{2\eta K}.
\tag{17}
\]

\item \textbf{Monotonicity of $\Phi(x_k)$.} The proximal gradient iteration is a descent method, i.e. $\Phi(x_{k+1})\le\Phi(x_k)$ (follows from (5) and the proximal optimality condition). Hence the largest function value among the first $K$ iterates is $\Phi(x_0)$ and the smallest is $\Phi(x_K)$. Consequently  

\[
\Phi(x_K)-\Phi^\star\le\frac{1}{K}\sum_{k=0}^{K-1}\bigl(\Phi(x_k)-\Phi^\star\bigr).
\tag{18}
\]

\item \textbf{Combine (17) and (18).} This yields the desired bound  

\[
\Phi(x_K)-\Phi^\star\le\frac{\|x_{0}-x^\star\|^2}{2\eta K},
\]
which completes the proof of Theorem~\ref{thm:conv}. ∎
\end{enumerate}

\section*{Rate Derivation (With Constants)}

From Theorem~\ref{thm:conv} we have explicitly  

\[
\boxed{\;\Phi(x_K)-\Phi^\star\le \frac{R^2}{2\eta K},\qquad 
R:=\|x_0-x^\star\|.\;}
\]

If we choose the maximal admissible stepsize $\eta=1/L$, the constant becomes $\frac{L R^2}{2K}$. Thus the convergence constant is proportional to the Lipschitz constant $L$ and the squared distance of the initial point to the solution set.

\section*{Special Cases}

\begin{itemize}
  \item \textbf{Pure Gradient Descent ($g\equiv0$).} The proximal operator reduces to the identity, and the bound \eqref{3} coincides with the classical $O(1/K)$ rate for gradient descent.
  \item \textbf{Indicator Function $g=\iota_{\mathcal{C}}$.} The method becomes projected gradient descent onto a convex set $\mathcal{C}$. The same $O(1/K)$ guarantee holds.
  \item \textbf{Strongly Convex $f+g$.} If $\Phi$ is $\mu$‑strongly convex, a refined analysis (adding a term $\mu\|x_k-x^\star\|^2$ in (9)) yields a linear rate $\mathcal{O}((1-\eta\mu)^K)$.
\end{itemize}

\end{document}